\documentclass[10pt,a4paper]{article}
\usepackage[utf8]{inputenc}
\usepackage[T5]{fontenc}
\usepackage[vietnamese,english]{babel}
\usepackage[margin=1.3cm,top=0.75cm,bottom=1.0cm,includehead,includefoot,headheight=18pt,headsep=5pt,footskip=16pt]{geometry}
\usepackage{amsmath,amssymb,array,tabularx,enumitem,fancyhdr,tikz,xcolor,needspace,multicol}
\usetikzlibrary{calc,angles,quotes,arrows.meta,babel,patterns}

\definecolor{ink}{RGB}{30,55,70}
\definecolor{accent}{RGB}{18,82,112}

\setlength{\parindent}{0pt}
\setlength{\parskip}{0pt}

\newcommand{\headerQuote}{%
  \ifcase\value{page}\relax
    \or Math gives us hope that every problem has a solution% Page 1
    \or Math + You = Champion% Page 2
    \or Where Frustration Becomes Inspiration% Page 3
    \or Turning Confusion into Clarity% Page 4
    \or Don't be an X, be a Y of mathematics% Page 5
    \or Don't get mad. Learn to add% Page 6
    \or The only way to learn mathematics is to do mathematics% Page 7
    \or Small steps every day lead to big results% Page 8
    \or Every expert was once a beginner% Page 9
    \or Success is the sum of small efforts repeated daily% Page 10
    \or Mathematics is the music of reason% Page 11
    \or Believe you can and you're halfway there% Page 12
    \or Math gives us hope that every problem has a solution% Page 13
    \or Math + You = Champion% Page 14
    \or Where Frustration Becomes Inspiration% Page 15
    \or Turning Confusion into Clarity% Page 16
    \or Don't be an X, be a Y of mathematics% Page 17
    \or Don't get mad. Learn to add% Page 18
    \or The only way to learn mathematics is to do mathematics% Page 19
    \or Small steps every day lead to big results% Page 20
    \or Every expert was once a beginner% Page 21
    \or Success is the sum of small efforts repeated daily% Page 22
    \or Mathematics is the music of reason% Page 23
    \or Believe you can and you're halfway there% Page 24
    \or Math gives us hope that every problem has a solution% Page 25
    \or Math + You = Champion% Page 26
    \or Where Frustration Becomes Inspiration% Page 27
    \or Turning Confusion into Clarity% Page 28
    \or Don't be an X, be a Y of mathematics% Page 29
    \or Don't get mad. Learn to add% Page 30
    \else Math gives us hope that every problem has a solution%
  \fi
}

\pagestyle{fancy}
\fancyhf{}
\lhead{\footnotesize\textit{\color{gray!70!black}\headerQuote}}
\rhead{\small\textbf{\color{ink}Compiled by: Dương Gia Bảo}}
\lfoot{}
\rfoot{\small Page \thepage/3}
\renewcommand{\headrulewidth}{0.4pt}
\renewcommand{\footrulewidth}{0pt}

\fancypagestyle{dapan}{
  \fancyhf{}
  \lhead{\footnotesize\textit{\color{gray!70!black}\headerQuote}}
  \rhead{\small\textbf{\color{ink}Compiled by: Dương Gia Bảo}}
  \lfoot{\small\textit{\color{gray!60!black}Answer Key — Grade 10 Math Review (Chapter IV: Vectors)}}
  \rfoot{\small Page \thepage/2}
  \renewcommand{\headrulewidth}{0.4pt}
  \renewcommand{\footrulewidth}{0pt}
}

\newcommand{\parttitle}[1]{\par\vspace{3.5pt}\needspace{3\baselineskip}{\large\bfseries\color{ink}#1}\par\nopagebreak\vspace{2pt}}
\newcommand{\four}[4]{\par\noindent\begin{tabularx}{\linewidth}{@{}*{4}{>{\raggedright\arraybackslash}X}@{}}\textbf{A.} #1&\textbf{B.} #2&\textbf{C.} #3&\textbf{D.} #4\end{tabularx}}
\newcommand{\two}[4]{\par\noindent\begin{tabularx}{\linewidth}{@{}*{2}{>{\raggedright\arraybackslash}X}@{}}\textbf{A.} #1&\textbf{B.} #2\\[3.8pt]\textbf{C.} #3&\textbf{D.} #4\end{tabularx}}
\newcommand{\one}[4]{\par\noindent\begin{tabularx}{\linewidth}{@{}>{\raggedright\arraybackslash}X@{}}\textbf{A.} #1\\[4pt]\textbf{B.} #2\\[4pt]\textbf{C.} #3\\[4pt]\textbf{D.} #4\end{tabularx}}
\newcommand{\answerbox}{\hfill\mbox{\textit{Answer:}\ \fbox{\rule{0pt}{7.5pt}\hspace{1.8cm}}}}

\setlist[enumerate,1]{label=\textbf{Question \arabic*.},leftmargin=2.1cm,labelsep=.15cm,itemsep=1.2pt,topsep=1.2pt,parsep=0pt}
\setlist[enumerate,2]{label=\alph*),leftmargin=.55cm,itemsep=1.5pt,topsep=0.2pt,parsep=0pt}
\newcolumntype{C}[1]{>{\centering\arraybackslash}m{#1}}


\begin{document}

\noindent
\begin{minipage}[c]{0.38\textwidth}
    \centering
    \textbf{DO MATH TO BE HAPPY}\\[-2pt]
    \rule{3.2cm}{0.6pt}\\[2pt]
    \textbf{OFFICIAL REVIEW TEST}
\end{minipage}\hfill
\begin{minipage}[c]{0.52\textwidth}
    \centering
    \textbf{GRADE 10 MATH REVIEW — CHAPTER IV. VECTORS}\\[1pt]
    \textbf{ACADEMIC YEAR 2026 -- 2027}\\[1pt]
    \textbf{Subject: MATHEMATICS -- Grade 10}\\
    \textit{\small Time allowed: 90 minutes (excluding distribution time)}
\end{minipage}

\vspace{0.05cm}
\noindent
\begin{minipage}[b]{0.78\textwidth}
    \textbf{Full Name:}\dotfill\textbf{ID:}\makebox[3.5cm]{\dotfill}
\end{minipage}\hfill
\begin{minipage}[b]{0.20\textwidth}
    \hfill
    \setlength{\fboxrule}{0.8pt}%
    \fbox{\begin{minipage}{2.3cm}
        \centering
        \textbf{\small Test Code}\\
        \textbf{1042-EN}
    \end{minipage}}
\end{minipage}
\vspace{0.05cm}

\parttitle{PART I. Multiple-Choice Questions}
\textit{Answer Questions 1 to 20. Choose only one correct option for each question.}
\begin{enumerate}
\item Given a regular hexagon $ABCDEF$ with center $O$. How many nonzero vectors have their initial and terminal points at the vertices of the regular hexagon?
\four{$36$.}{$15$.}{$30$.}{$12$.}

\item Given a parallelogram $ABCD$ with center $O$. Which of the following vectors is equal to $\vec{AB}$?
\four{$\vec{DC}$.}{$\vec{CD}$.}{$\vec{BA}$.}{$\vec{OC}$.}

\item For any four points $A, B, C, D$, which of the following statements is always true?
\two{$\vec{AB} + \vec{CD} = \vec{AC} + \vec{DB}$.}{$\vec{AB} + \vec{BC} = \vec{CA}$.}{$\vec{AB} - \vec{AC} = \vec{BC}$.}{$\vec{AB} + \vec{BC} + \vec{CD} = \vec{AD}$.}

\item Let $ABC$ be an equilateral triangle with side length $a$. What is the magnitude of the vector $\vec{u} = \vec{AB} + \vec{AC}$?
\four{$a$.}{$a\sqrt{3}$.}{$2a$.}{$\dfrac{a\sqrt{3}}{2}$.}

\item Let $ABCD$ be a square with side length $a$. What is the magnitude of the vector $\vec{v} = \vec{AB} - \vec{AD}$?
\four{$a\sqrt{2}$.}{$2a$.}{$a$.}{$\dfrac{a\sqrt{2}}{2}$.}

\item Let $G$ be the centroid of triangle $ABC$ and $M$ be the midpoint of side $BC$. Which of the following vector equations is correct?
\two{$\vec{GA} + \vec{GB} + \vec{GC} = 3\vec{0}$.}{$\vec{AG} = \dfrac{1}{2}\vec{AM}$.}{$\vec{AG} = 2\vec{GM}$.}{$\vec{GB} + \vec{GC} = 2\vec{GA}$.}

\item Given triangle $ABC$. Let $M$ be a point on segment $BC$ such that $MB = 2MC$. Which of the following identities is correct?
\two{$\vec{AM} = \dfrac{2}{3}\vec{AB} + \dfrac{1}{3}\vec{AC}$.}{$\vec{AM} = \dfrac{1}{3}\vec{AB} + \dfrac{2}{3}\vec{AC}$.}{$\vec{AM} = \dfrac{1}{2}\vec{AB} + \dfrac{1}{2}\vec{AC}$.}{$\vec{AM} = \dfrac{1}{3}\vec{AB} - \dfrac{2}{3}\vec{AC}$.}

\item Given three distinct points $A, B, C$. What is the necessary and sufficient condition for $A, B, C$ to be collinear?
\two{$\vec{AB} = \vec{AC}$.}{$|\vec{AB}| = |\vec{AC}|$.}{$\vec{AB} + \vec{AC} = \vec{0}$.}{There exists a real number $k \ne 0$ such that $\vec{AB} = k\vec{AC}$.}

\item In the Cartesian coordinate plane $Oxy$, given the vector $\vec{u} = 2\vec{i} - 5\vec{j}$. What are the coordinates of $\vec{u}$?
\four{$(-5, 2)$.}{$(2, 5)$.}{$(2, -5)$.}{$(-2, 5)$.}

\item In the coordinate plane $Oxy$, given two points $A(-1, 3)$ and $B(4, 1)$. What are the coordinates of vector $\vec{AB}$?
\four{$(3, 4)$.}{$(5, -2)$.}{$(-5, 2)$.}{$(3, -2)$.}

\item In the coordinate plane $Oxy$, given two points $A(2, 4)$ and $B(-4, 2)$. What are the coordinates of the midpoint $M$ of segment $AB$?
\four{$M(-1, 3)$.}{$M(3, 3)$.}{$M(-2, 6)$.}{$M(1, -3)$.}

\item In the coordinate plane $Oxy$, given three points $A(1, 2)$, $B(3, 5)$, and $C(5, 2)$. What are the coordinates of point $D$ such that quadrilateral $ABCD$ is a parallelogram?
\four{$D(7, 5)$.}{$D(-1, 5)$.}{$D(3, 1)$.}{$D(3, -1)$.}
\pagebreak

\item In the coordinate plane $Oxy$, given two vectors $\vec{a} = (2, -3)$ and $\vec{b} = (m, 6)$. Vectors $\vec{a}$ and $\vec{b}$ are parallel (collinear) if and only if
\four{$m = 4$.}{$m = -4$.}{$m = -9$.}{$m = 9$.}

\item Let $ABC$ be an equilateral triangle. What is the angle between the vectors $\vec{AB}$ and $\vec{BC}$?
\four{$60^\circ$.}{$30^\circ$.}{$120^\circ$.}{$150^\circ$.}

\item Let $ABCD$ be a square with side length $a$. What is the dot product $\vec{AB} \cdot \vec{AC}$?
\four{$a^2$.}{$a^2\sqrt{2}$.}{$\dfrac{a^2\sqrt{2}}{2}$.}{$2a^2$.}

\item In the coordinate plane $Oxy$, given two vectors $\vec{u} = (3, -2)$ and $\vec{v} = (4, 5)$. What is the value of the dot product $\vec{u} \cdot \vec{v}$?
\four{$22$.}{$-2$.}{$14$.}{$2$.}

\item In the coordinate plane $Oxy$, given point $A(-3, 4)$. What is the magnitude of vector $\vec{OA}$?
\four{$7$.}{$5$.}{$1$.}{$\sqrt{7}$.}

\item In the coordinate plane $Oxy$, given two vectors $\vec{x} = (2, m)$ and $\vec{y} = (3, -6)$. Vectors $\vec{x}$ and $\vec{y}$ are orthogonal (perpendicular) if and only if
\four{$m = -1$.}{$m = 4$.}{$m = 1$.}{$m = 0$.}

\item In the coordinate plane $Oxy$, given two vectors $\vec{a} = (1, 2)$ and $\vec{b} = (-2, 1)$. What is the angle between $\vec{a}$ and $\vec{b}$?
\four{$90^\circ$.}{$0^\circ$.}{$45^\circ$.}{$180^\circ$.}

\item A constant force $\vec{F}$ of magnitude $50\text{ N}$ acts on an object, moving it a distance of $10\text{ m}$ in a direction making an angle of $60^\circ$ with $\vec{F}$. What is the work $W$ done by the force $\vec{F}$?
\four{$500\text{ J}$.}{$250\sqrt{3}\text{ J}$.}{$500\sqrt{3}\text{ J}$.}{$250\text{ J}$.}
\end{enumerate}

\vspace{3pt}
\parttitle{PART II. True/False Questions}
\textit{Answer Questions 1 to 4. For each statement a), b), c), d) in every question, indicate whether it is True or False.}
\begin{enumerate}
\item Given a rhombus $ABCD$ with center $O$, side length $a$, and angle $\widehat{BAD} = 60^\circ$ (as shown in the figure).
\par\noindent\begin{minipage}[t]{0.72\linewidth}
\begin{enumerate}
    \item Triangle $ABD$ is an equilateral triangle with side length $a$.
    \item The magnitude of the vector $\vec{u} = \vec{AB} + \vec{AD}$ is $a\sqrt{2}$.
    \item The magnitude of the vector $\vec{v} = \vec{AB} - \vec{AD}$ is $a$.
    \item The vector equality $\vec{AO} + \vec{CO} = \vec{AC}$ is true.
\end{enumerate}
\end{minipage}\hfill
\begin{minipage}[t]{0.26\linewidth}\raggedleft\vspace{-6pt}
\begin{tikzpicture}[scale=0.55, >=Stealth]
  \coordinate (A) at (0,0);
  \coordinate (B) at (2, 1.155);
  \coordinate (C) at (4,0);
  \coordinate (D) at (2, -1.155);
  \coordinate (O) at (2,0);
  \draw[thick, ink] (A) -- (B) -- (C) -- (D) -- cycle;
  \draw[dashed, gray!70] (A) -- (C);
  \draw[dashed, gray!70] (B) -- (D);
  \node[left, font=\scriptsize] at (A) {$A$};
  \node[above, font=\scriptsize] at (B) {$B$};
  \node[right, font=\scriptsize] at (C) {$C$};
  \node[below, font=\scriptsize] at (D) {$D$};
  \node[above right, font=\tiny] at (O) {$O$};
  \pic[draw, angle radius=0.35cm, "$60^\circ$", angle eccentricity=1.7, font=\tiny] {angle = D--A--B};
\end{tikzpicture}
\end{minipage}

\item Given triangle $ABC$. Let $M$ be the midpoint of side $AB$, $N$ be a point on side $AC$ such that $AN = 2NC$, and $G$ be the centroid of triangle $ABC$.
\begin{enumerate}
    \item Vector $\vec{CM}$ can be expressed as $\vec{CM} = \dfrac{1}{2}\vec{CA} - \vec{CB}$.
    \item Vector $\vec{BN}$ can be expressed as $\vec{BN} = \dfrac{2}{3}\vec{AC} - \vec{AB}$.
    \item Vector $\vec{AG}$ can be expressed as $\vec{AG} = \dfrac{1}{3}\vec{AB} + \dfrac{1}{3}\vec{AC}$.
    \item The three points $M, G, C$ are not collinear.
\end{enumerate}

\item In the coordinate plane $Oxy$, given three points $A(-2, 1)$, $B(2, 4)$, and $C(5, 0)$.
\begin{enumerate}
    \item The coordinates of the vectors are $\vec{AB} = (4, 3)$ and $\vec{AC} = (7, -1)$.
    \item The three points $A, B, C$ are non-collinear and form a triangle.
    \item The centroid $G$ of triangle $ABC$ has coordinates $G\left(\dfrac{5}{3}, \dfrac{7}{3}\right)$.
    \item The dot product is $\vec{AB} \cdot \vec{AC} = 25$.
\end{enumerate}

\item Let $ABC$ be an isosceles right triangle at $A$ with $AB = AC = 4$. Let $M$ be the midpoint of side $BC$.
\begin{enumerate}
    \item The dot product is $\vec{BA} \cdot \vec{BC} = 8$.
    \item The magnitude of vector $\vec{AM}$ is $4$.
    \item The dot product is $\vec{AM} \cdot \vec{BC} = 0$.
    \item The magnitude of the vector $\vec{u} = \vec{AB} + 2\vec{AC}$ is $4\sqrt{5}$.
\end{enumerate}
\end{enumerate}

\pagebreak
\parttitle{PART III. Short-Answer Questions}
\textit{Answer Questions 1 to 6. Write your final answer in the corresponding box.}
\begin{enumerate}
\item Given rectangle $ABCD$ with $AB = 6$ and $AD = 8$. Find the magnitude of $\vec{u} = \vec{AB} + \vec{AD}$.\answerbox

\item Given triangle $ABC$. Let $I$ be a point on side $BC$ such that $IB = 3IC$. Given that $\vec{AI} = x\vec{AB} + y\vec{AC}$ where $x, y \in \mathbb{R}$, find the value of $P = 4x + 8y$.\answerbox

\item In the coordinate plane $Oxy$, given three points $A(1, -2)$, $B(3, 2)$, and $C(m, 6)$. Find the value of the parameter $m$ such that $A, B, C$ are collinear.\answerbox

\item In the coordinate plane $Oxy$, given two vectors $\vec{a} = (2, 1)$ and $\vec{b} = (3, -1)$. Find the angle between $\vec{a}$ and $\vec{b}$ (in degrees).\answerbox

\item In the coordinate plane $Oxy$, given three points $A(2, 4)$, $B(1, 1)$, and $C(k, 0)$. Find the value of $k$ so that triangle $ABC$ is right-angled at $B$.\answerbox

\item \begin{minipage}[t]{0.70\linewidth}
Two tugboats simultaneously pull a barge at point $O$. The first tugboat exerts a pulling force $\vec{F}_1$ of magnitude $6000\text{ N}$, and the second tugboat exerts a pulling force $\vec{F}_2$ of magnitude $8000\text{ N}$. Knowing that $\vec{F}_1$ and $\vec{F}_2$ are perpendicular to each other ($\vec{F}_1 \perp \vec{F}_2$, as shown in the figure), calculate the magnitude of the resultant force $\vec{F} = \vec{F}_1 + \vec{F}_2$ acting on the barge (in Newtons).\answerbox
\end{minipage}\hfill
\begin{minipage}[t]{0.28\linewidth}\raggedleft\vspace{-6pt}
\begin{tikzpicture}[scale=0.32, >=Stealth]
  \coordinate (O) at (0,0);
  \coordinate (F1) at (6,0);
  \coordinate (F2) at (0,8);
  \coordinate (F) at (6,8);
  \draw[->, very thick, accent] (O) -- (F1) node[below, font=\scriptsize] {$\vec{F}_1$};
  \draw[->, very thick, accent] (O) -- (F2) node[left, font=\scriptsize] {$\vec{F}_2$};
  \draw[dashed, gray!60] (F1) -- (F) -- (F2);
  \draw[->, line width=1.1pt, ink] (O) -- (F) node[above right, font=\scriptsize] {$\vec{F}$};
  \draw[gray!70] (0.6,0) -- (0.6,0.6) -- (0,0.6);
  \node[below left, font=\scriptsize] at (O) {$O$};
\end{tikzpicture}
\end{minipage}
\end{enumerate}

\vspace{3pt}
\parttitle{PART IV. Free-Response Problems}
\textit{Show complete reasoning and step-by-step solutions for the following problems:}

\vspace{4pt}
\textbf{Problem 1 (1.0 point).} Given triangle $ABC$. Let $M$ be the midpoint of side $AB$, $N$ be a point on side $AC$ such that $AN = \dfrac{1}{3}AC$, and $P$ be a point on line $BC$ such that $B$ is the midpoint of segment $PC$.
\begin{enumerate}[label=\alph*), leftmargin=0.6cm, itemsep=2pt]
    \item (0.5 point) Express vector $\vec{MN}$ in terms of $\vec{AB}$ and $\vec{AC}$.
    \item (0.5 point) Express vector $\vec{MP}$ in terms of $\vec{AB}$ and $\vec{AC}$, and hence prove that points $M, N, P$ are collinear.
\end{enumerate}

\vspace{6pt}
\textbf{Problem 2 (1.0 point).} In the coordinate plane $Oxy$, given triangle $ABC$ with vertices $A(-1, 2)$, $B(3, 4)$, and $C(2, -1)$.
\begin{enumerate}[label=\alph*), leftmargin=0.6cm, itemsep=2pt]
    \item (0.5 point) Calculate the perimeter of triangle $ABC$ (round your result to the nearest tenth).
    \item (0.5 point) Find the coordinates of the orthocenter $H$ of triangle $ABC$.
\end{enumerate}

\vspace{6pt}
\textbf{Problem 3 (1.0 point).} An object of mass $m = 10\text{ kg}$ (with gravitational weight $\vec{P}$ of magnitude $P = 100\text{ N}$) is suspended in equilibrium by two light, inextensible strings $OA$ and $OB$ attached to a horizontal ceiling. The two strings make angles of $30^\circ$ and $60^\circ$ with the ceiling, respectively (as shown in the figure).
\par\noindent\begin{minipage}[t]{0.70\linewidth}
\begin{enumerate}[label=\alph*), leftmargin=0.6cm, itemsep=2pt]
    \item (0.5 point) State the equilibrium condition at the suspension point $O$ as a vector equation relating the weight $\vec{P}$ and the tension forces $\vec{T}_1, \vec{T}_2$.
    \item (0.5 point) Calculate the magnitudes of the tension forces $T_1$ in string $OA$ and $T_2$ in string $OB$ (give exact values or round to one decimal place).
\end{enumerate}
\end{minipage}\hfill
\begin{minipage}[t]{0.28\linewidth}\raggedleft\vspace{-15pt}
\begin{tikzpicture}[scale=0.46, >=Stealth]
  \fill[gray!20] (-1.8,2.0) rectangle (3.2,2.2);
  \draw[thick] (-1.8,2.0) -- (3.2,2.0);
  \coordinate (A) at (-0.6, 2.0);
  \coordinate (B) at (1.8, 2.0);
  \coordinate (O) at (0.6, 0.4);
  \draw[thick, ink] (A) -- (O) node[midway, left, font=\tiny] {$T_1$};
  \draw[thick, ink] (B) -- (O) node[midway, right, font=\tiny] {$T_2$};
  \draw[->, very thick, accent] (O) -- (0.6, -0.6) node[right, font=\tiny] {$\vec{P}$};
  \fill[ink] (O) circle (1.8pt) node[below left, font=\tiny] {$O$};
  \fill[accent!80!black] (0.38, -0.6) rectangle (0.82, -1.0);
  \node[font=\tiny, white] at (0.6, -0.8) {$m$};
  \draw[dashed, gray!70] (A) -- +(0,-0.5);
  \draw[dashed, gray!70] (B) -- +(0,-0.5);
  \node[below, font=\tiny] at (-0.2, 2.0) {$30^\circ$};
  \node[below, font=\tiny] at (1.4, 2.0) {$60^\circ$};
\end{tikzpicture}
\end{minipage}


\newpage
\pagestyle{dapan}
\setcounter{page}{1}

\begin{center}
  {\Large\bfseries\color{ink} COMPREHENSIVE ANSWER KEY}\\[2pt]
  {\small\color{gray!75!black} Grade 10 Math Review — Chapter IV. Vectors (Code 1042-EN)}
\end{center}
\vspace{1pt}

\noindent\textbf{PART I. Multiple-Choice Questions (20 questions)}
\vspace{1.5pt}

\begin{center}
\renewcommand{\arraystretch}{1.15}
\begin{tabular}{|C{0.95cm}|C{0.95cm}|C{0.95cm}|C{0.95cm}|C{0.95cm}|C{0.95cm}|C{0.95cm}|C{0.95cm}|C{0.95cm}|C{0.95cm}|}
\hline
\textbf{1} & \textbf{2} & \textbf{3} & \textbf{4} & \textbf{5} & \textbf{6} & \textbf{7} & \textbf{8} & \textbf{9} & \textbf{10} \\ \hline
\textbf{C} & \textbf{A} & \textbf{D} & \textbf{B} & \textbf{A} & \textbf{C} & \textbf{B} & \textbf{D} & \textbf{C} & \textbf{B} \\ \hline
\hline
\textbf{11} & \textbf{12} & \textbf{13} & \textbf{14} & \textbf{15} & \textbf{16} & \textbf{17} & \textbf{18} & \textbf{19} & \textbf{20} \\ \hline
\textbf{A} & \textbf{D} & \textbf{B} & \textbf{C} & \textbf{A} & \textbf{D} & \textbf{B} & \textbf{C} & \textbf{A} & \textbf{D} \\ \hline
\end{tabular}
\end{center}

\vspace{1pt}
\noindent\begin{minipage}[t]{0.58\linewidth}
\textbf{PART II. True/False Questions}\\[1.5pt]
\renewcommand{\arraystretch}{1.15}
\begin{tabular}{|c|c|c|c|c|}
\hline
\textbf{Q} & \textbf{a)} & \textbf{b)} & \textbf{c)} & \textbf{d)} \\ \hline
\textbf{1} & \textbf{T} & \textbf{F} & \textbf{T} & \textbf{F} \\ \hline
\textbf{2} & \textbf{F} & \textbf{T} & \textbf{T} & \textbf{F} \\ \hline
\textbf{3} & \textbf{T} & \textbf{T} & \textbf{F} & \textbf{T} \\ \hline
\textbf{4} & \textbf{F} & \textbf{F} & \textbf{T} & \textbf{T} \\ \hline
\end{tabular}
\end{minipage}\hfill
\begin{minipage}[t]{0.40\linewidth}
\textbf{PART III. Short-Answer Questions}\\[1.5pt]
\renewcommand{\arraystretch}{1.15}
\begin{tabular}{|c|c|c|c|}
\hline
\textbf{Q} & \textbf{Answer} & \textbf{Q} & \textbf{Answer} \\ \hline
\textbf{1} & $10$ & \textbf{4} & $45$ \\ \hline
\textbf{2} & $7$  & \textbf{5} & $4$ \\ \hline
\textbf{3} & $5$  & \textbf{6} & $10000$ \\ \hline
\end{tabular}
\end{minipage}

\vspace{5pt}
\noindent\textbf{PART IV. Free-Response Summary:}\\[2pt]
- \textbf{Problem 1:} a) $\vec{MN} = -\dfrac{1}{2}\vec{AB} + \dfrac{1}{3}\vec{AC}$. b) $\vec{MP} = \dfrac{3}{2}\vec{AB} - \vec{AC} = -3\vec{MN} \Rightarrow M, N, P$ are collinear.\\
- \textbf{Problem 2:} a) Perimeter of $\triangle ABC$: $2\sqrt{5} + \sqrt{26} + 3\sqrt{2} \approx 13.8$. b) Orthocenter $H\left(\dfrac{2}{3}, \dfrac{5}{3}\right)$.\\
- \textbf{Problem 3:} a) $\vec{T}_1 + \vec{T}_2 + \vec{P} = \vec{0}$. b) Tensions: $T_1 = 50\text{ N}$; $T_2 = 50\sqrt{3}\text{ N} \approx 86.6\text{ N}$.

\vspace{4pt}
\begin{center}
  {\large\bfseries\color{ink} DETAILED SOLUTIONS}
\end{center}
\vspace{1pt}

\begin{multicols}{2}
\raggedcolumns
\small

\textbf{\color{accent}PART I. MULTIPLE-CHOICE}

\textbf{Question 1.} The hexagon has $6$ vertices. The number of nonzero vectors formed by choosing an ordered pair of vertices is $P(6,2) = 6 \times 5 = 30$. Choice \textbf{C}.

\textbf{Question 2.} In parallelogram $ABCD$, opposite sides $AB$ and $DC$ are parallel and equal in length, and the direction from $A$ to $B$ matches $D$ to $C$, so $\vec{AB} = \vec{DC}$. Choice \textbf{A}.

\textbf{Question 3.} By the extended triangle rule (head-to-tail addition): $\vec{AB} + \vec{BC} + \vec{CD} = \vec{AC} + \vec{CD} = \vec{AD}$. Choice \textbf{D}.

\textbf{Question 4.} Let $M$ be the midpoint of $BC$; then $\vec{AB} + \vec{AC} = 2\vec{AM}$. Since $\triangle ABC$ is equilateral of side $a$, median $AM = \dfrac{a\sqrt{3}}{2} \Rightarrow |\vec{u}| = 2 \cdot \dfrac{a\sqrt{3}}{2} = a\sqrt{3}$. Choice \textbf{B}.

\textbf{Question 5.} We have $\vec{AB} - \vec{AD} = \vec{DB}$. Since $ABCD$ is a square of side $a$, the diagonal length is $DB = a\sqrt{2}$. Choice \textbf{A}.

\textbf{Question 6.} Since $G$ is the centroid and $M$ is the midpoint of $BC$, $AG = \dfrac{2}{3}AM \Rightarrow AG = 2GM$. Since $\vec{AG}$ and $\vec{GM}$ have the same direction, $\vec{AG} = 2\vec{GM}$. Choice \textbf{C}.

\textbf{Question 7.} Since $M$ lies on segment $BC$ with $MB = 2MC$, we have $\vec{MB} = -2\vec{MC} \Leftrightarrow \vec{AB} - \vec{AM} = -2(\vec{AC} - \vec{AM}) \Leftrightarrow 3\vec{AM} = \vec{AB} + 2\vec{AC} \Leftrightarrow \vec{AM} = \dfrac{1}{3}\vec{AB} + \dfrac{2}{3}\vec{AC}$. Choice \textbf{B}.

\textbf{Question 8.} Three distinct points $A, B, C$ are collinear if and only if $\vec{AB}$ and $\vec{AC}$ are parallel, i.e., there exists a real scalar $k \ne 0$ such that $\vec{AB} = k\vec{AC}$. Choice \textbf{D}.

\textbf{Question 9.} Vector $\vec{u} = 2\vec{i} - 5\vec{j}$ has components $(2, -5)$. Choice \textbf{C}.

\textbf{Question 10.} $\vec{AB} = (4 - (-1), 1 - 3) = (5, -2)$. Choice \textbf{B}.

\textbf{Question 11.} Midpoint $M$ of $AB$: $x_M = \dfrac{2 + (-4)}{2} = -1$; $y_M = \dfrac{4 + 2}{2} = 3 \Rightarrow M(-1, 3)$. Choice \textbf{A}.

\textbf{Question 12.} $ABCD$ is a parallelogram iff $\vec{AD} = \vec{BC} \Leftrightarrow (x_D - 1, y_D - 2) = (5 - 3, 2 - 5) = (2, -3) \Rightarrow x_D = 3, y_D = -1 \Rightarrow D(3, -1)$. Choice \textbf{D}.

\textbf{Question 13.} Two vectors $\vec{a}$ and $\vec{b}$ are parallel iff $2(6) - (-3)m = 0 \Leftrightarrow 12 + 3m = 0 \Leftrightarrow m = -4$. Choice \textbf{B}.

\textbf{Question 14.} Since $\triangle ABC$ is equilateral, $\widehat{B} = 60^\circ$. The angle between $\vec{AB}$ and $\vec{BC}$ is supplementary to $\widehat{B}$: $(\vec{AB}, \vec{BC}) = 180^\circ - 60^\circ = 120^\circ$. Choice \textbf{C}.

\textbf{Question 15.} The square's diagonal is $AC = a\sqrt{2}$ and $\widehat{BAC} = 45^\circ$. Hence $\vec{AB} \cdot \vec{AC} = AB \cdot AC \cdot \cos 45^\circ = a \cdot a\sqrt{2} \cdot \dfrac{\sqrt{2}}{2} = a^2$. Choice \textbf{A}.

\textbf{Question 16.} $\vec{u} \cdot \vec{v} = 3(4) + (-2)(5) = 12 - 10 = 2$. Choice \textbf{D}.

\textbf{Question 17.} $|\vec{OA}| = \sqrt{(-3)^2 + 4^2} = \sqrt{25} = 5$. Choice \textbf{B}.

\textbf{Question 18.} $\vec{x} \perp \vec{y} \Leftrightarrow 2(3) + m(-6) = 0 \Leftrightarrow 6 - 6m = 0 \Leftrightarrow m = 1$. Choice \textbf{C}.

\textbf{Question 19.} We have $\vec{a} \cdot \vec{b} = 1(-2) + 2(1) = 0 \Rightarrow \vec{a} \perp \vec{b} \Rightarrow (\vec{a}, \vec{b}) = 90^\circ$. Choice \textbf{A}.

\textbf{Question 20.} $W = |\vec{F}| \cdot d \cdot \cos 60^\circ = 50 \cdot 10 \cdot 0.5 = 250\text{ J}$. Choice \textbf{D}.

\columnbreak

\textbf{\color{accent}PART II. TRUE / FALSE}

\textbf{Question 1.}
a) Rhombus $ABCD$ has $AB = AD = a$ and $\widehat{BAD} = 60^\circ$, so $\triangle ABD$ is equilateral of side $a \Rightarrow$ \textbf{True}.\\
b) $\vec{AB} + \vec{AD} = \vec{AC}$. Since $\triangle ABD$ is equilateral, $BD = a$, altitude $AO = \dfrac{a\sqrt{3}}{2} \Rightarrow AC = 2AO = a\sqrt{3} \ne a\sqrt{2} \Rightarrow$ \textbf{False}.\\
c) $\vec{AB} - \vec{AD} = \vec{DB} \Rightarrow |\vec{v}| = BD = a \Rightarrow$ \textbf{True}.\\
d) Since $O$ is the midpoint of $AC$, $\vec{AO} + \vec{CO} = \vec{0} \ne \vec{AC} \Rightarrow$ \textbf{False}.

\textbf{Question 2.}
a) Since $M$ is the midpoint of $AB$, $\vec{CM} = \dfrac{1}{2}(\vec{CA} + \vec{CB}) \ne \dfrac{1}{2}\vec{CA} - \vec{CB} \Rightarrow$ \textbf{False}.\\
b) $\vec{BN} = \vec{AN} - \vec{AB} = \dfrac{2}{3}\vec{AC} - \vec{AB} \Rightarrow$ \textbf{True}.\\
c) For centroid $G$ of $\triangle ABC$, $\vec{AG} = \dfrac{1}{3}(\vec{AB} + \vec{AC}) = \dfrac{1}{3}\vec{AB} + \dfrac{1}{3}\vec{AC} \Rightarrow$ \textbf{True}.\\
d) Because $M$ is the midpoint of $AB$, $CM$ is a median of $\triangle ABC$. Centroid $G$ always lies on median $CM$, so $M, G, C$ are collinear $\Rightarrow$ \textbf{False}.

\textbf{Question 3.}
a) $\vec{AB} = (2 - (-2), 4 - 1) = (4, 3)$; $\vec{AC} = (5 - (-2), 0 - 1) = (7, -1) \Rightarrow$ \textbf{True}.\\
b) Since $\dfrac{4}{7} \ne \dfrac{3}{-1}$, the two vectors are not parallel, so $A, B, C$ form a triangle $\Rightarrow$ \textbf{True}.\\
c) Centroid: $x_G = \dfrac{-2+2+5}{3} = \dfrac{5}{3}$; $y_G = \dfrac{1+4+0}{3} = \dfrac{5}{3} \Rightarrow G\left(\dfrac{5}{3}, \dfrac{5}{3}\right) \ne G\left(\dfrac{5}{3}, \dfrac{7}{3}\right) \Rightarrow$ \textbf{False}.\\
d) $\vec{AB} \cdot \vec{AC} = 4(7) + 3(-1) = 28 - 3 = 25 \Rightarrow$ \textbf{True}.

\textbf{Question 4.}
a) $\triangle ABC$ is right isosceles at $A \Rightarrow \widehat{B} = 45^\circ$, $BC = 4\sqrt{2}$. Thus $\vec{BA} \cdot \vec{BC} = BA \cdot BC \cdot \cos 45^\circ = 4 \cdot 4\sqrt{2} \cdot \dfrac{\sqrt{2}}{2} = 16 \ne 8 \Rightarrow$ \textbf{False}.\\
b) Since $\triangle ABC$ is right at $A$, median $AM = \dfrac{BC}{2} = 2\sqrt{2} \ne 4 \Rightarrow$ \textbf{False}.\\
c) Since $\triangle ABC$ is isosceles at $A$, median $AM$ is also the altitude, so $AM \perp BC \Rightarrow \vec{AM} \cdot \vec{BC} = 0 \Rightarrow$ \textbf{True}.\\
d) $|\vec{u}|^2 = \vec{AB}^2 + 4\vec{AB}\cdot\vec{AC} + 4\vec{AC}^2 = 16 + 0 + 4(16) = 80 \Rightarrow |\vec{u}| = \sqrt{80} = 4\sqrt{5} \Rightarrow$ \textbf{True}.

\textbf{\color{accent}PART III. SHORT-ANSWER}

\textbf{Question 1.} $\vec{AB} + \vec{AD} = \vec{AC}$. Thus $|\vec{u}| = AC = \sqrt{AB^2 + AD^2} = \sqrt{6^2 + 8^2} = \mathbf{10}$.

\textbf{Question 2.} $\vec{IB} = -3\vec{IC} \Leftrightarrow \vec{AB} - \vec{AI} = -3(\vec{AC} - \vec{AI}) \Leftrightarrow 4\vec{AI} = \vec{AB} + 3\vec{AC} \Leftrightarrow \vec{AI} = \dfrac{1}{4}\vec{AB} + \dfrac{3}{4}\vec{AC}$. Hence $x = \dfrac{1}{4}, y = \dfrac{3}{4} \Rightarrow P = 4\left(\dfrac{1}{4}\right) + 8\left(\dfrac{3}{4}\right) = 1 + 6 = \mathbf{7}$.

\textbf{Question 3.} $\vec{AB} = (2, 4)$, $\vec{AC} = (m - 1, 8)$. Points are collinear iff $\dfrac{m - 1}{2} = \dfrac{8}{4} = 2 \Rightarrow m - 1 = 4 \Rightarrow m = \mathbf{5}$.

\textbf{Question 4.} $\vec{a} \cdot \vec{b} = 2(3) + 1(-1) = 5$. Also $|\vec{a}| = \sqrt{4+1} = \sqrt{5}$, $|\vec{b}| = \sqrt{9+1} = \sqrt{10}$. Thus $\cos(\vec{a}, \vec{b}) = \dfrac{5}{\sqrt{5}\sqrt{10}} = \dfrac{1}{\sqrt{2}} = \dfrac{\sqrt{2}}{2} \Rightarrow (\vec{a}, \vec{b}) = \mathbf{45}^\circ$.

\textbf{Question 5.} $\vec{BA} = (1, 3)$, $\vec{BC} = (k - 1, -1)$. Triangle is right at $B \Leftrightarrow \vec{BA} \cdot \vec{BC} = 0 \Leftrightarrow 1(k - 1) + 3(-1) = 0 \Leftrightarrow k - 4 = 0 \Rightarrow k = \mathbf{4}$.

\textbf{Question 6.} Since $\vec{F}_1 \perp \vec{F}_2$, the resultant force magnitude is $|\vec{F}| = \sqrt{|\vec{F}_1|^2 + |\vec{F}_2|^2} = \sqrt{6000^2 + 8000^2} = \mathbf{10000}\text{ N}$.

\textbf{\color{accent}PART IV. FREE-RESPONSE}

\textbf{Problem 1.}
a) We have $\vec{MN} = \vec{AN} - \vec{AM}$.\\
Since $M$ is the midpoint of $AB \Rightarrow \vec{AM} = \dfrac{1}{2}\vec{AB}$.\\
Since $AN = \dfrac{1}{3}AC \Rightarrow \vec{AN} = \dfrac{1}{3}\vec{AC}$.\\
Therefore, $\vec{MN} = -\dfrac{1}{2}\vec{AB} + \dfrac{1}{3}\vec{AC}$.\\
b) Since $B$ is the midpoint of $PC$, $\vec{PB} = \vec{BC} = \vec{AC} - \vec{AB}$.\\
Thus $\vec{AP} = \vec{AB} + \vec{BP} = \vec{AB} - \vec{PB} = \vec{AB} - (\vec{AC} - \vec{AB}) = 2\vec{AB} - \vec{AC}$.\\
Hence $\vec{MP} = \vec{AP} - \vec{AM} = (2\vec{AB} - \vec{AC}) - \dfrac{1}{2}\vec{AB} = \dfrac{3}{2}\vec{AB} - \vec{AC}$.\\
Notice that $\vec{MP} = -3\left(-\dfrac{1}{2}\vec{AB} + \dfrac{1}{3}\vec{AC}\right) = -3\vec{MN}$.\\
Since $\vec{MP}$ and $\vec{MN}$ are parallel, the three points $M, N, P$ are collinear.

\textbf{Problem 2.}
a) Side lengths:\\
$AB = \sqrt{(3 - (-1))^2 + (4 - 2)^2} = \sqrt{16 + 4} = \sqrt{20} = 2\sqrt{5} \approx 4.47$.\\
$BC = \sqrt{(2 - 3)^2 + (-1 - 4)^2} = \sqrt{1 + 25} = \sqrt{26} \approx 5.10$.\\
$AC = \sqrt{(2 - (-1))^2 + (-1 - 2)^2} = \sqrt{9 + 9} = \sqrt{18} = 3\sqrt{2} \approx 4.24$.\\
Perimeter of $\triangle ABC$: $P = 2\sqrt{5} + \sqrt{26} + 3\sqrt{2} \approx \mathbf{13.8}$.\\
b) Let the orthocenter be $H(x, y)$. We have $\vec{AH} = (x + 1, y - 2)$, $\vec{BC} = (-1, -5)$, and $\vec{BH} = (x - 3, y - 4)$, $\vec{AC} = (3, -3)$.\\
$H$ is the orthocenter of $\triangle ABC \Leftrightarrow \begin{cases} \vec{AH} \cdot \vec{BC} = 0 \\ \vec{BH} \cdot \vec{AC} = 0 \end{cases}$\\
$\Leftrightarrow \begin{cases} -1(x + 1) - 5(y - 2) = 0 \\ 3(x - 3) - 3(y - 4) = 0 \end{cases} \Leftrightarrow \begin{cases} -x - 5y + 9 = 0 \\ 3x - 3y + 3 = 0 \end{cases}$\\
$\Leftrightarrow \begin{cases} x + 5y = 9 \\ x - y = -1 \end{cases} \Leftrightarrow \begin{cases} x = \dfrac{2}{3} \\ y = \dfrac{5}{3} \end{cases}$.\\
Thus, the orthocenter of $\triangle ABC$ is $H\left(\mathbf{\dfrac{2}{3}}, \mathbf{\dfrac{5}{3}}\right)$.

\textbf{Problem 3.}
a) The forces acting on suspension point $O$ are the gravitational weight $\vec{P}$ and the string tensions $\vec{T}_1, \vec{T}_2$.\\
For the system to be in equilibrium, the net force on $O$ must be the zero vector:
$$\vec{T}_1 + \vec{T}_2 + \vec{P} = \vec{0}.$$
b) Choose a Cartesian coordinate system $Oxy$ with the $x$-axis horizontal to the right and the $y$-axis vertical upward.\\
We have $\vec{P} = (0, -100)$ (since $P = 100\text{ N}$).\\
Since the strings make angles of $30^\circ$ and $60^\circ$ with the ceiling:\\
$\vec{T}_1 = (-T_1 \cos 30^\circ,\ T_1 \sin 30^\circ) = \left(-\dfrac{\sqrt{3}}{2}T_1,\ \dfrac{1}{2}T_1\right)$,\\
$\vec{T}_2 = (T_2 \cos 60^\circ,\ T_2 \sin 60^\circ) = \left(\dfrac{1}{2}T_2,\ \dfrac{\sqrt{3}}{2}T_2\right)$.\\
From the equilibrium condition $\vec{T}_1 + \vec{T}_2 = -\vec{P} = (0, 100)$, we get the system:\\
$\begin{cases} -\dfrac{\sqrt{3}}{2}T_1 + \dfrac{1}{2}T_2 = 0 \\[4pt] \dfrac{1}{2}T_1 + \dfrac{\sqrt{3}}{2}T_2 = 100 \end{cases} \Leftrightarrow \begin{cases} T_2 = T_1\sqrt{3} \\[4pt] \dfrac{1}{2}T_1 + \dfrac{\sqrt{3}}{2}(T_1\sqrt{3}) = 100 \end{cases}$\\
$\Leftrightarrow 2T_1 = 100 \Leftrightarrow \mathbf{T_1 = 50\text{ N}}$.\\
Hence $\mathbf{T_2 = 50\sqrt{3}\text{ N}} \approx \mathbf{86.6\text{ N}}$.

\end{multicols}
\end{document}
